% ─── Introduction ────────────────────────────────────────────────────────────
% OIES style: set out the problem, why it matters, the precise gap this paper
% fills, and what the contribution is. Evidence-based, policy-relevant.
% House style: British spelling, Oxford / New Hart's Rules, minimal dashes.

Which technology should decarbonise Europe's homes, and where, is now a commitment of
capital that cannot easily be reversed. Between now and 2050 hundreds of billions of
euros of heating infrastructure will be replaced, and the choice between electrifying
the heat and reusing the gas network through hydrogen decides how that money is spent.
The stakes are large because the sector is large. Buildings account for approximately 40
per cent of final energy consumption in the European Union, the greater part of it
burned for space heating and hot water, so the sector sits at the centre of climate
policy from the European Green Deal through to the Energy Performance of Buildings
Directive \citep{IEA2022buildings,EuropeanCommission2020GreenDeal}.

What makes the choice hard is the form the heat takes. It is delivered through hundreds
of millions of individual dwellings, each with its own vintage, fabric, climate and
occupant, and replaced only slowly as boilers fail and buildings change hands
\citep{Azevedo2022, Chatterjee2024Navigating}. Residential heat must therefore be abated
one household at a time, against real constraints of cost, disruption and capital
turnover \citep{Maia2024Household, Arbulu2022Barriers}, and it is the heterogeneity of
the stock that makes the technology choice consequential, because a route that is cheap
in one dwelling can be costly in the next.

European policy has converged on heat pumps as the principal route. The
\emph{Fit for 55} package, the revised Energy Performance of Buildings Directive
and the REPowerEU plan all treat electrification as the default for buildings,
and the case is strong on the merits \citep{REpowerEU2022,EPBD2024}.
A heat pump delivers a coefficient of performance (COP) of two to four, so it
undercuts any combustion appliance on delivered energy, and it decarbonises
automatically as the grid does \citep{IEA2022heatpumps,Abbasi2021}.
Deployment is nonetheless slowed by frictions that fall unevenly across the
stock, among them the cost of retrofitting poorly insulated dwellings, distribution-grid
reinforcement in dense urban areas, a COP that degrades in
the coldest weather and in the oldest fabric, split incentives between landlords
and tenants, and uneven public acceptance
\citep{Lingard2020Residential,Gordon2023Divergent}. These
frictions do not overturn electrification's cost case; they create pockets,
poorly insulated buildings, grid-constrained zones and regions of dense
industrial heat demand, where the upfront cost or engineering burden of
electrification is locally steep, and into which a hydrogen alternative has been
proposed.

Hydrogen is the alternative most often advanced, and the most contested. Its
appeal is institutional as much as technical. A low-carbon gas could in principle
reuse the existing distribution network and the appliance habit of a gas-heated
home, sidestepping the fabric upgrades and grid reinforcement that electrification
demands, and serving precisely the hard-to-electrify segment that heat pumps
struggle to reach \citep{Dodds2015,Bothe2021}. REPowerEU
raised the EU's 2030 hydrogen ambition to 20~Mt of renewable and imported supply,
with buildings named among the candidate demand centres
\citep{REpowerEU2022, EUCommissionSWD2022}. The energy-security shock that
followed the 2022 gas crisis sharpened the appetite for a domestic, storable,
network-compatible fuel \citep{Goldthau2023Eu, Zakeri2022Pandemic}.

A substantial sceptical literature has answered that hydrogen for space heating is
thermodynamically wasteful and expensive, that it competes for scarce electrolytic
supply with industry and transport, where substitutes are fewer, and that the
network incumbency argument understates the cost of conversion and the value of
the heat a pump extracts from the air
\citep{Rosenow2022PipeDream,Cebon2022}.
This dispute matters in practice. It shapes whether gas distribution
networks are written down or repurposed, where renovation subsidy is directed,
and how the heating infrastructure replaced before 2050 is committed. Switzerland,
ringed by states with active hydrogen strategies and without domestic gas
resources of its own, faces a sharper version of the same question, namely whether
imported low-carbon hydrogen has any role in a national heat system that it would
otherwise electrify \citep{Gabrielli2020,BfE2023}.

Despite the stakes, the quantitative basis for resolving the dispute is thinner
than the policy debate implies, and the literature can be mapped onto three
recurrent limitations, each exemplified by influential work. The first limitation
is one of spatial resolution. The continental scenario assessments that frame the
heating debate operate at national or whole-system resolution, averaging over
exactly the spatial heterogeneity, in building stock, climate, demand density and
storage geology, that determines where, if anywhere, hydrogen can compete. The
pan-European race between hydrogen and heat pumps in \citet{Billerbeck2023Race}
and the whole-system electrification accounts of \citet{Hoseinpoori2022} and
\citet{Brown2018Synergies} resolve the technology contest finely but the geography
coarsely, while the spatially explicit studies that do exist remain single-country,
as in the Spanish optimisation of \citet{Maestre2022} and the German
techno-economic comparison of \citet{Huckebrink2022}. The cross-country pattern of
where hydrogen's economics diverge from the average is thus left unmapped.

The
second limitation concerns how hydrogen's heating share enters the analysis.
Studies typically fix that share by assumption, whether as a policy target, a
scenario lever or a feasibility ceiling, and then evaluate the consequences, a
procedure visible in the scenario-driven shares of \citet{Billerbeck2023Race} and
the technology-portfolio framing of \citet{Weidner2023Planetary}. That procedure
cannot answer the prior question of whether a hydrogen heating market would form at
all, because it imposes the very outcome it ought to derive.

Pricing hydrogen on its levelised cost over a year, as most such comparisons do,
compounds the problem by misstating how a dispatchable fuel competes. A system operator
does not run the cheapest technology on average, but the one with the lowest short-run
marginal cost at each hour, so a fuel that loses the annual cost contest may still win,
or lose, the few hundred hours of the winter peak on which dispatchable capacity
actually earns its keep \citep{Joskow2008capacity,Quarton2019Value}. The third
limitation is a severed link between end uses.

The electrification that displaces fossil heat does not make the heat-supply problem
disappear; it relocates part of it into the power sector, which must then meet a
cold-snap heating load on the same scarce winter evenings. That interaction has largely
been treated apart, the buildings-side studies of the heating choice and the power-side
studies of system adequacy addressing it separately. The heat-peak work of
\citet{Zeyen2021Mitigating} and the Swiss winter-deficit analysis of
\citet{Mellot2024Mitigating} take the two halves on their own, not as one merit order.
The updated REPowerEU context, the accelerated hydrogen targets and the introduction of
a second emissions trading system for buildings and road transport (ETS2) further
postdate much of the published modelling.

These three limitations define the questions this paper sets out to answer. The core
question is European and country-level: under what conditions, on economic and
environmental cost grounds, is decarbonised hydrogen competitive in the buildings sector
across the EU-27, the United Kingdom and Switzerland under a net-zero by 2050 case, and where those conditions hold. It resolves into three research questions, each closing
one of the limitations above, restated as the three gaps of Section~\ref{sec:lit:gap},
and each answered in full by the analysis that follows.

RQ1 (EU, UK and Switzerland, country-level) takes up the first two gaps together,
testing the hydrogen-versus-heat-pump trade-off on a consistent operating-cost basis and
deriving hydrogen's heating share from where the fuel can win the cold-snap peak instead
of imposing it as a target \citep{Rosenow2022PipeDream,ICCT2023h2heating}. RQ2
(Switzerland, national case) closes the third gap for Switzerland as a high-cost,
hydro-rich, cavern-poor system reliant on imported hydrogen
\citep{Gabrielli2020,IEA2023switzerland}. RQ3 (cross-sector synergy, scoped to the power
sector) names the buildings-to-power coupling that the severed link implies, the one
synergy this paper carries through in full
\citep{Brown2018Synergies,Zeyen2021Mitigating}. The three questions are stated in full,
with the scope boundary that separates what this paper models from what it leaves to
subsequent work, in Section~\ref{sec:rq}.

\paragraph{Three contributions.}
This paper resolves the three limitations just set out, the missing spatial
resolution, the imposed and not derived heating share, and the severed
heating-to-power link, within a single, internally consistent framework, and in
doing so makes three contributions.
\begin{enumerate}
  \item \emph{Methodological.} We build a bottom-up
  techno-economic model of residential heating across the EU-27, the United Kingdom
  and Switzerland, 29 countries, over 2025 to 2050.
  Useful heat demand is reconstructed at NUTS3 resolution from EUBUCCO building footprints and TABULA
  heat intensities, validated against the Hotmaps regional database, and propagated
  through eight heating technologies and four multi-lever scenarios by a 200-draw
  Monte Carlo that carries parametric uncertainty end to end, alongside an
  independent least-cost optimisation \citep{MilojevicDupont2023Eubucco,Loga2016Tabula}. The NUTS3 reconstruction earns its place. It grounds the demand and its climate-and-density gradient in observed building stock, not a top-down national split, the step that makes the country aggregates trustworthy.
  The economic verdict is then formed at the country level, where hydrogen's prospects
  diverge from a continental average because the conditions its case turns on, salt-cavern
  storage geology, the carbon-priced gas it must undercut and dense district-heating
  loads, differ markedly from one country to the next, a divergence that a single
  continental average dilutes to invisibility. This matters more for hydrogen than for
  heat pumps, because a heat pump is broadly available wherever the grid reaches and its
  cost varies smoothly with climate and fabric, so an average represents it tolerably.
  \item \emph{Reframing.} Rather than fixing a heating share by policy assumption,
  we \emph{derive} it through cost and dispatch tests that are run separately. On the full annual cost of delivered heat, with each carrier charged its own equipment
capital, its own last-mile network and the seasonal storage its own load shape requires,
the best heat pump is cheaper than a gas boiler in 24 of the 29 countries and cheaper
than a hydrogen boiler in 22, so most of the stock electrifies. Two of the seven
exceptions are geological and five turn on the retail tariff the heat pump pays. Electrified heat is then an electricity load, which moves part of the
heating problem onto the power sector and its scarcest winter evenings. That displaced
load is priced the way a system operator actually prices electricity, by short-run
marginal cost at each hour. Hydrogen enters at
  both the first and the last of those steps, as a base-load competitor on annual
  cost in the countries where retail electricity is dearest, and again as a
  candidate peaker competing on operating cost against the carbon-taxed gas plant.
  Hydrogen is represented along its full supply chain, from production route and
  import corridor through seasonal storage to the distribution network, with a
  geographically bounded storage adder of roughly \euro{}50/MWh in countries with
  salt-cavern endowment and roughly \euro{}100/MWh where caverns are absent.
  \item \emph{Heating-to-power bridge.} We carry the analysis across the
  heating-to-power bridge that the electrification pathway creates, applying the
  same dispatch logic in three linked arenas, the building winter peak, the
  district-heating stack, and the power sector on which the electrified heat load
  now depends. It is this bridge that makes the power-sector merit order a heating
  question at all, and it is the synergy named in RQ3.
\end{enumerate}

Three findings follow, set out with their uncertainty bounds in
Sections~\ref{sec:results} and~\ref{sec:power} and previewed here.

First, the base-load and dispatch tests converge instead of cutting against each
other. On the full annual cost of delivered heat in 2050, with each carrier charged
its own last-mile network and with hydrogen charged the seasonal store its
winter-weighted load shape requires, the EU median is roughly \euro{}118.5/MWh for the
best heat pump, \euro{}122.8/MWh for a hydrogen boiler and \euro{}132.7/MWh for a gas
boiler. The best heat pump is the cheaper of the two clean options in 22 of 29
countries at a median country gap of \euro{}15.4/MWh, and hydrogen holds the remaining
seven, five of them salt-cavern markets.

Priced instead on the short-run marginal cost a system operator actually uses, hydrogen undercuts the carbon-taxed fossil competitor
at the peak hour in a cavern-endowed minority, and the count rises with policy
ambition across the four scenarios. Charging both carriers alike, so that hydrogen pays
the same last-mile adder its base-load comparison uses and the district-heat operator's
gas is bought at the wholesale hub price, it wins at the building winter peak in 0
countries under Current Policies, 2 under Stated Policies, 7 under Net Zero and 14 under
H2~Push; in the district-heating stack in 0, 0, 7 and 7 countries respectively; and
in the power sector in 0, 0, 7 and 15 on the evolving-efficiency turbine series, or
0, 0, 7 and 7 on the even-handed series that holds both turbines at 0.40. On the
asymmetric accounting an earlier draft carried, which charges hydrogen no last-mile
network and prices the district-heat gas unit at the residential retail tariff, the first
two arenas read 0, 4, 9 and 16 and 0, 5, 11 and 20.

These counts measure where hydrogen would actually be dispatched, not where it is
deployed by policy or where it breaks even on capital. The power-sector count of 15 is
bimodal and should be read with care, the decisive operating-cost advantage holding in
only the seven salt-cavern markets and the remaining eight sitting at a marginal
\euro{}0 to 8/MWh-e (Section~\ref{sec:power:meritorder}). Both tests therefore turn on
one physical endowment, which is also why their agreement corroborates less than two
separate confirmations would.

What lead hydrogen keeps in those markets is a fiscal result before a
technological one, since by 2050 the median European household pays roughly
\euro{}209/MWh for taxed retail electricity against roughly \euro{}50/MWh for hydrogen
delivered and untaxed, and equalising the two treatments would widen the heat pump's
margin further.

Second, the electrified load opens a power-sector niche that no market would build
unaided. Once heat pumps supply roughly half of useful heat, residential heating
becomes about 9 per cent of 2050 electricity demand and a driver of the system's
winter peak, calling for a firm peaking fleet of about 278~GW against a study-area
system peak of about 889~GW, roughly 31 per cent of peak capacity, which is a capacity
role and not an energy one. Both figures are computed on a single synthetic weather
year, so they size that role and do not assess adequacy.

The same merit-order test inside the power sector puts the hydrogen turbine level with
the carbon-priced gas peaker at the cross-country median,
about \euro{}238 against \euro{}239 per megawatt-hour of electricity, and decisively
below it in the cavern markets. Yet the peaker runs too few hours to recover its
capital from market rent, the classic missing-money problem, recovering about 9 per
cent across the cavern markets on a capacity-weighted basis, so even the dispatched
role would require a capacity payment or equivalent support. An independent least-cost
optimisation, which minimises total system cost over the whole technology portfolio
with no hour-by-hour dispatch, reaches a 2050 mix of about 53 per cent district heat
and 45 per cent heat pumps, with residual fossil at about 1.4 per cent. Its hydrogen
share of about a third of one per cent is bounded by a per-country ceiling that is zero
in 27 markets, so that program corroborates the district-heat ordering and not the
hydrogen verdict.

Third, policy sets the emissions outcome. Across four
scenarios sharing a common $\approx$644~MtCO$_2$ 2025 baseline, 2050
residential-heating emissions range from 352 to 10~MtCO$_2$. Read together, the
base-load, dispatch and investment tests converge on one condition. Hydrogen holds
base load in seven countries, five of which have the salt caverns that make seasonal
storage cheap, wins the cold-snap dispatch only in the cavern endowment itself, and
nowhere recovers a peaker's capital from market rent, including there. The three tests
therefore identify the narrow and policy-supported conditions under which a hydrogen
heating market could exist.

The remainder of the paper is organised as follows. Section~\ref{sec:rq} states the
research questions and fixes the scope. Section~\ref{sec:lit} reviews the literature
on buildings decarbonisation, heat pumps, hydrogen for heat, merit order and system
value, spatial bottom-up methods, and the Switzerland context, and from it derives
the research gap. Section~\ref{sec:method} sets out the data sources, the demand trajectory and the
four-scenario lever design. Section~\ref{sec:results} reports demand, the technology
mix, levelised costs, emissions and the base-load hydrogen gap, and
Section~\ref{sec:power} carries the analysis across the heating-to-power bridge, from
the merit order to the firm-capacity requirement and the peaker's capital recovery.
Section~\ref{sec:discussion} interprets the results and Section~\ref{sec:conclusion}
concludes with the policy implications. Four appendices carry the rest, and the first of
them is substantial rather than supplementary. Appendix~\ref{app:detail} holds the
bottom-up demand build and its validation, the merit-order, missing-money and emissions
formulations, the Monte Carlo machinery, the hourly dispatch and unit commitment, the
hydrogen supply and last-mile network costs, the district-heating stack, the
infrastructure bill and the Switzerland case; a reader following the method will spend
as much time there as in Section~\ref{sec:method}. Appendix~\ref{sec:costopt} presents
the least-cost optimisation and the implied-carbon-price finding,
Appendix~\ref{sec:sensitivity} tests robustness, and Appendix~\ref{sec:limitations}
states the principal limitations.
